We now restrict projective evaluation to the affine syntactic sector.  Every
evaluated history then has the form
\[
  x\longmapsto \Phi x+\xi,
  \qquad \Phi\in K^\times,\quad \xi\in K.
\]
The accumulated scale $\Phi$ and translation $\xi$ retain more of the process
than an endpoint value.  Their composition law forces two natural translation
coordinates.  The target-frame coordinate $\xi$ weights each translation by
the scales that occur after it; the source-normalized coordinate
$\widehat\xi=\Phi^{-1}\xi$ discounts it by the scales accumulated up to that
point.

Throughout Subsections~\ref{subsec:discrete-affine-cocycles}--
\ref{subsec:relative-affine-defect}, $K$ is an arbitrary field and every
affine multiplier is nonzero.  The differential calculation in
Subsection~\ref{subsec:affine-maurer-cartan} is over $\mathbb R$.

\subsection{Discrete affine evaluation}
\label{subsec:discrete-affine-cocycles}

Let
\[
  f_i(x)=s_i x+t_i,
  \qquad s_i\in K^\times,\quad t_i\in K,
\]
and regard $(f_1,\ldots,f_n)$ as a chronological history: $f_1$ acts first.
Write
\begin{equation}
  f_n\circ\cdots\circ f_1(x)=\Phi_nx+\xi_n.
  \label{eq:affine-history-normal-form}
\end{equation}
For the empty history we set $\Phi_0=1$ and $\xi_0=0$.

\begin{proposition}[Affine cocycle formulas]
\label{prop:affine-cocycle-formulas}
\label{prop:target-affine-cocycle}
\label{prop:source-normalized-cocycle}
For a chronological affine history $(f_1,\ldots,f_n)$,
\begin{align}
  \Phi_n&=\prod_{i=1}^{n}s_i,
  \label{eq:accumulated-affine-scale}\\
  \xi_n&=\sum_{i=1}^{n}t_i\prod_{j=i+1}^{n}s_j.
  \label{eq:target-affine-cocycle}
\end{align}
The empty product is $1$.  Defining the source-normalized translation by
\begin{equation}
  \widehat\xi_n=\Phi_n^{-1}\xi_n,
  \label{eq:source-normalized-translation-definition}
\end{equation}
one also has
\begin{equation}
  \widehat\xi_n
  =\sum_{i=1}^{n}\frac{t_i}{\prod_{j=1}^{i}s_j}.
  \label{eq:source-normalized-affine-cocycle}
\end{equation}
\end{proposition}

\begin{proof}
We prove the first two identities simultaneously by induction on $n$.  They
hold for $n=0$.  Suppose
\[
  f_n\circ\cdots\circ f_1(x)=\Phi_nx+\xi_n.
\]
Postcomposing chronologically by $f_{n+1}$ gives
\[
  f_{n+1}(\Phi_nx+\xi_n)
  =s_{n+1}\Phi_nx+(s_{n+1}\xi_n+t_{n+1}).
\]
Therefore
\[
  \Phi_{n+1}=s_{n+1}\Phi_n
  =\prod_{i=1}^{n+1}s_i.
\]
Using the induction hypothesis for $\xi_n$,
\begin{align*}
  \xi_{n+1}
  &=s_{n+1}\sum_{i=1}^{n}t_i
       \prod_{j=i+1}^{n}s_j+t_{n+1}\\
  &=\sum_{i=1}^{n}t_i
       \prod_{j=i+1}^{n+1}s_j
    +t_{n+1}\prod_{j=n+2}^{n+1}s_j,
\end{align*}
which is exactly \eqref{eq:target-affine-cocycle} for $n+1$.

Every $s_i$ is nonzero, so $\Phi_n$ is invertible.  Dividing the $i$th term of
\eqref{eq:target-affine-cocycle} by
$\Phi_n=\prod_{j=1}^{n}s_j$ gives
\[
  \frac{t_i\prod_{j=i+1}^{n}s_j}
       {\prod_{j=1}^{n}s_j}
  =\frac{t_i}{\prod_{j=1}^{i}s_j}.
\]
Summation proves \eqref{eq:source-normalized-affine-cocycle}.
\end{proof}

Formula \eqref{eq:target-affine-cocycle} is a future-weighting rule.  A
translation $t_i$ is expressed in target coordinates only after every later
scale $s_{i+1},\ldots,s_n$ has acted.  Formula
\eqref{eq:source-normalized-affine-cocycle} is a past-normalization rule: the
same translation is transported back through the scale accumulated from the
source through step $i$.  Neither formula is an additional geometric
postulate.

\begin{example}[Three steps]
For
\[
  f_1(x)=s_1x+t_1,
  \quad f_2(x)=s_2x+t_2,
  \quad f_3(x)=s_3x+t_3,
\]
direct substitution gives
\[
  f_3f_2f_1(x)
  =s_3s_2s_1x+s_3s_2t_1+s_3t_2+t_3.
\]
Here and below juxtaposition of the displayed functions means ordinary
functional composition $f_3\circ f_2\circ f_1$, not chronological word
concatenation.  The formula makes the future scales attached to each
translation visible.
\end{example}

\subsection{Cocycle laws under chronological concatenation}

For an affine history $\gamma$, write
\[
  \rho(\gamma)(x)=\Phi_\gamma x+\xi_\gamma,
  \qquad
  \widehat\xi_\gamma=\Phi_\gamma^{-1}\xi_\gamma.
\]
Recall that $\gamma\delta$ means first $\gamma$ and then $\delta$, so
$\rho(\gamma\delta)=\rho(\delta)\rho(\gamma)$.

\begin{proposition}[Chronological affine cocycle law]
\label{prop:chronological-affine-cocycle-law}
For affine histories $\gamma$ and $\delta$,
\begin{align}
  \Phi_{\gamma\delta}
    &=\Phi_\delta\Phi_\gamma,
  \label{eq:scale-concatenation-law}\\
  \xi_{\gamma\delta}
    &=\Phi_\delta\xi_\gamma+\xi_\delta,
  \label{eq:target-concatenation-law}\\
  \widehat\xi_{\gamma\delta}
    &=\widehat\xi_\gamma
      +\Phi_\gamma^{-1}\widehat\xi_\delta.
  \label{eq:source-concatenation-law}
\end{align}
If $\gamma$ is invertible, then
\begin{equation}
  \Phi_{\gamma^{-1}}=\Phi_\gamma^{-1},
  \qquad
  \xi_{\gamma^{-1}}=-\Phi_\gamma^{-1}\xi_\gamma
  =-\widehat\xi_\gamma.
  \label{eq:affine-inverse-coordinates}
\end{equation}
\end{proposition}

\begin{proof}
Apply $\rho(\delta)$ to $\rho(\gamma)(x)$:
\[
  \rho(\delta)(\Phi_\gamma x+\xi_\gamma)
  =\Phi_\delta\Phi_\gamma x
   +\Phi_\delta\xi_\gamma+\xi_\delta.
\]
This proves \eqref{eq:scale-concatenation-law} and
\eqref{eq:target-concatenation-law}.  Since $K$ is commutative,
\begin{align*}
  \widehat\xi_{\gamma\delta}
  &=(\Phi_\delta\Phi_\gamma)^{-1}
    (\Phi_\delta\xi_\gamma+\xi_\delta)\\
  &=\Phi_\gamma^{-1}\xi_\gamma
    +\Phi_\gamma^{-1}\Phi_\delta^{-1}\xi_\delta,
\end{align*}
which is \eqref{eq:source-concatenation-law}.  Finally, solving
$y=\Phi_\gamma x+\xi_\gamma$ for $x$ yields
\[
  \rho(\gamma)^{-1}(y)
  =\Phi_\gamma^{-1}y-\Phi_\gamma^{-1}\xi_\gamma,
\]
and hence \eqref{eq:affine-inverse-coordinates}.
\end{proof}

These formulas show why the translation coordinate is a cocycle rather than
an additive homomorphism.  The semidirect-product action of scale on
translation is part of the arithmetic composition law.

\subsection{Left and right affine differentials}
\label{subsec:affine-maurer-cartan}

We now pass to the positive real affine group and write
\begin{equation}
  g(M,\xi)
  =\begin{pmatrix}e^M&\xi\\0&1\end{pmatrix}.
  \label{eq:positive-affine-matrix-coordinates}
\end{equation}
The symbol $M$ denotes accumulated logarithmic scale; it is distinct from the
fixed infinitesimal intensity $\lambda$ used in later flow equations.

\begin{proposition}[Left and right affine differentials]
\label{prop:left-right-affine-differentials}
For the matrix $g$ in \eqref{eq:positive-affine-matrix-coordinates}, the
translation components of the left and right Maurer--Cartan forms are
\begin{align}
  \bigl(g^{-1}dg\bigr)_{\mathrm{tr}}
    &=e^{-M}\,d\xi,
  \label{eq:left-maurer-cartan-translation}\\
  \bigl(dg\,g^{-1}\bigr)_{\mathrm{tr}}
    &=d\xi-\xi\,dM.
  \label{eq:right-maurer-cartan-translation}
\end{align}
If $\widehat\xi=e^{-M}\xi$, then
\begin{equation}
  d\xi-\xi\,dM=e^M d\widehat\xi.
  \label{eq:maurer-cartan-coordinate-conversion}
\end{equation}
\end{proposition}

\begin{proof}
We have
\[
  g^{-1}
  =\begin{pmatrix}
      e^{-M}&-e^{-M}\xi\\0&1
    \end{pmatrix},
  \qquad
  dg
  =\begin{pmatrix}
      e^M dM&d\xi\\0&0
    \end{pmatrix}.
\]
Multiplication on the two sides gives
\[
  g^{-1}dg
  =\begin{pmatrix}
      dM&e^{-M}d\xi\\0&0
    \end{pmatrix}
\]
and
\[
  dg\,g^{-1}
  =\begin{pmatrix}
      dM&d\xi-\xi dM\\0&0
    \end{pmatrix}.
\]
The upper-right entries prove
\eqref{eq:left-maurer-cartan-translation} and
\eqref{eq:right-maurer-cartan-translation}.  Differentiating
$\widehat\xi=e^{-M}\xi$ gives
\[
  d\widehat\xi=e^{-M}(d\xi-\xi dM),
\]
which is equivalent to
\eqref{eq:maurer-cartan-coordinate-conversion}.
\end{proof}

The left Maurer--Cartan form $g^{-1}dg$ measures velocity in body or source
coordinates; the right form $dg\,g^{-1}$ measures it in spatial or target
coordinates.  This left--right group distinction is not the same as
operand-slot chirality.  Relating a slot choice to left or right group
multiplication requires an additional rule specifying how a new context is
inserted into a represented history.  No such identification is assumed here.

There are likewise two separate coordinate statements.  The matrix coordinate
$\xi$ is the target-frame translation in $x\mapsto e^M x+\xi$, whereas
$\widehat\xi=e^{-M}\xi$ is its source-normalized value.  Equations
\eqref{eq:left-maurer-cartan-translation}--
\eqref{eq:maurer-cartan-coordinate-conversion} relate their differentials but
do not make the two coordinates interchangeable.

\subsection{Relative affine defect}
\label{subsec:relative-affine-defect}

The affine cocycle permits a comparison of two histories before any global
torsion theory is introduced.

\begin{definition}[Relative affine defect]
\label{def:relative-affine-defect}
For two affine histories $\gamma$ and $\delta$, the \emph{relative affine
defect of $\gamma$ with reference $\delta$} is
\begin{equation}
  \mathcal K_{\mathrm{aff}}(\gamma,\delta)
  =\rho(\delta)^{-1}\rho(\gamma)
  \in\operatorname{Aff}(1,K).
  \label{eq:relative-affine-defect}
\end{equation}
The pair is \emph{scale-compatible} when
$\Phi_\gamma=\Phi_\delta$.
\end{definition}

The order of the arguments in
\eqref{eq:relative-affine-defect} is part of the convention.  Interchanging
$\gamma$ and $\delta$ inverts the group-valued defect.

\begin{proposition}[Coordinates of the relative affine defect]
\label{prop:relative-affine-defect-coordinates}
For arbitrary affine histories $\gamma$ and $\delta$,
\begin{equation}
  \mathcal K_{\mathrm{aff}}(\gamma,\delta)(x)
  =\frac{\Phi_\gamma}{\Phi_\delta}x
   +\frac{\xi_\gamma-\xi_\delta}{\Phi_\delta}.
  \label{eq:relative-affine-defect-coordinates}
\end{equation}
If the pair is scale-compatible with common scale $\Phi$, then
\begin{align}
  \nu_x(\gamma)-\nu_x(\delta)
    &=\xi_\gamma-\xi_\delta,
  \label{eq:scale-compatible-endpoint-defect}\\
  \mathcal K_{\mathrm{aff}}(\gamma,\delta)(x)
    &=x+\frac{\xi_\gamma-\xi_\delta}{\Phi}.
  \label{eq:scale-compatible-relative-translation}
\end{align}
In particular, the endpoint difference is then independent of $x$.
\end{proposition}

\begin{proof}
The inverse of $y\mapsto\Phi_\delta y+\xi_\delta$ is
$y\mapsto\Phi_\delta^{-1}(y-\xi_\delta)$.  Applying it after
$x\mapsto\Phi_\gamma x+\xi_\gamma$ gives
\eqref{eq:relative-affine-defect-coordinates}.  If
$\Phi_\gamma=\Phi_\delta=\Phi$, direct subtraction of the two endpoint formulas
gives \eqref{eq:scale-compatible-endpoint-defect}, and substituting the common
scale into \eqref{eq:relative-affine-defect-coordinates} gives
\eqref{eq:scale-compatible-relative-translation}.
\end{proof}

For unequal scales, the endpoint difference contains
$(\Phi_\gamma-\Phi_\delta)x$ and depends on the initial value.  In that case the
group-valued comparison remains defined, but there is no base-point-independent
scalar translation defect of the form
\eqref{eq:scale-compatible-endpoint-defect}.

\begin{example}[The elementary affine order defect]
\label{ex:elementary-affine-order-defect}
Fix $\mu\in K$ and $q\in K^\times$.  Let $\gamma$ first add $\mu$ and then
scale by $q$, while $\delta$ first scales by $q$ and then adds $\mu$.  Then
\[
  \rho(\gamma)(x)=q(x+\mu)=qx+q\mu,
  \qquad
  \rho(\delta)(x)=qx+\mu.
\]
The histories have the same scale $q$, and hence
\begin{equation}
  \nu_x(\gamma)-\nu_x(\delta)=\mu(q-1),
  \qquad
  \mathcal K_{\mathrm{aff}}(\gamma,\delta)(x)
  =x+\mu(1-q^{-1}).
  \label{eq:elementary-affine-order-defect}
\end{equation}
Over the positive real group, writing $q=e^M$ gives the endpoint defect
$\mu(e^M-1)$.
\end{example}

Equation \eqref{eq:elementary-affine-order-defect} is an open comparison of two
chronological histories.  It is not a closed commutator-loop holonomy, a
comparison with path inverse, or a mirror invariant.  Section~\ref{sec:acs-torsion}
will place compatible two-history defects in the global torsion framework.

\subsection{Propagation of a change in the initial accumulator}

The accumulated scale also controls sensitivity to the initial value.

\begin{proposition}[Exact affine perturbation propagation]
\label{prop:affine-perturbation-propagation}
Let $\gamma$ be an affine history with scale $\Phi_\gamma$.  For any
$x,\widetilde x\in K$,
\begin{equation}
  \nu_{\widetilde x}(\gamma)-\nu_x(\gamma)
  =\Phi_\gamma(\widetilde x-x).
  \label{eq:exact-affine-perturbation}
\end{equation}
When $\widetilde x\ne x$,
\begin{equation}
  \frac{\nu_{\widetilde x}(\gamma)-\nu_x(\gamma)}
       {\widetilde x-x}
  =\Phi_\gamma.
  \label{eq:affine-perturbation-ratio}
\end{equation}
\end{proposition}

\begin{proof}
The affine normal form gives
\[
  \nu_{\widetilde x}(\gamma)=\Phi_\gamma\widetilde x+\xi_\gamma,
  \qquad
  \nu_x(\gamma)=\Phi_\gamma x+\xi_\gamma.
\]
Subtracting cancels the translation cocycle and proves
\eqref{eq:exact-affine-perturbation}.  Division by the nonzero element
$\widetilde x-x$ proves \eqref{eq:affine-perturbation-ratio}.
\end{proof}

This is an exact discrete identity, not a linear approximation.  The
continuous affine flow of the next section is obtained by differentiating the
same semidirect-product action.  Supplementary expansions and convention
checks appear in Appendix~\ref{app:affine-cocycles}.
